\section{预期的成果}

\subsection{成果一（对应研究点一）}
\begin{enumerate}
  \item \textbf{AgentDojo框架扩展。} 完成AgentDojo与HuggingFace Transformers库的集成，实现开源模型（Qwen3-4B、Gemma-3-4B）的白盒梯度访问支持，开发标准化攻击集成接口，支持模型特定聊天模板与工具调用格式处理，为后续攻击研究提供可复用的基础设施。
  \item \textbf{GCG攻击适配与优化。} 实现GCG攻击在间接提示注入场景的完整适配，包括面向工具调用的目标函数设计、tokenization验证机制、前缀/后缀/混合注入结构策略、单任务与通用优化变体。预期在开源模型上通用GCG攻击实现显著攻击成功率。
  \item \textbf{GCG消融实验结果。} 通过梯度引导vs随机搜索消融实验，量化梯度信号在智能体攻击场景中的实际贡献；通过注入结构消融实验，揭示不同注入位置对攻击效果的影响规律；通过通用攻击泛化实验，评估联合优化发现的注入模式的跨场景适用性。
\end{enumerate}

\subsection{成果二（对应研究点二）}
\begin{enumerate}
  \item \textbf{TAP攻击适配与优化。} 实现TAP攻击在间接提示注入场景的完整适配，包括面向工具调用的树搜索算法改造、可靠性测试机制、攻击者/评估器LLM系统提示优化、并行评估实现。预期在开源模型上通用TAP攻击实现显著攻击成功率，且显著超越白盒GCG方法。
  \item \textbf{跨模型迁移性分析。} 在前沿闭源模型上评估直接优化攻击与跨模型迁移攻击的效果，揭示开源模型优化攻击的迁移局限性及前沿模型与开源模型之间的安全差距。
  \item \textbf{LLM-as-judge准确性评估。} 量化不同目标模型下评估器LLM的判断准确性，分析评估噪声对优化效率的影响，提出针对性改进策略。
  \item \textbf{关键因素分析。} 系统揭示影响攻击成功的核心因素：注入结构与语义内容、初始化随机性、攻击者模型能力、评估器可靠性、目标模型特性等，为防御设计提供攻击侧的实证依据。
\end{enumerate}

\subsection{原型应用系统}
基于整体研究内容，设计并实现一套面向LLM智能体的自动化提示注入攻击评估系统，系统包含五大核心模块：
\begin{enumerate}
  \item 框架集成与环境管理模块：封装AgentDojo环境接口，集成HuggingFace Transformers库，支持多模型（Qwen3-4B、Gemma-3-4B、GPT-5）的统一加载与工具调用格式适配。
  \item GCG白盒攻击模块：实现单任务GCG、通用GCG攻击优化，支持前缀/后缀/混合注入结构配置，集成tokenization验证和梯度计算管理。
  \item TAP黑盒攻击模块：实现单任务TAP、通用TAP攻击优化，支持攻击者LLM变异策略配置、评估器LLM评分与剪枝、可靠性测试和并行评估。
  \item 评估与基准测试模块：基于AgentDojo确定性检查函数，实现攻击成功率（ASR）、Security-at-N、任务域细分统计等指标的计算与可视化。
  \item 消融与迁移性分析模块：支持注入结构消融、梯度vs随机消融、跨模型迁移测试、跨域泛化测试等实验配置的自动化运行与结果汇总。
\end{enumerate}
原型系统支持从攻击配置、优化运行、结果评估到因素分析的全流程自动化，能够通过实验验证所提方法的有效性，最终形成可复现、可扩展的LLM智能体安全评估工具。

\subsection{学术论文}
基于研究内容与实验结果，计划撰写并发表会议论文一篇。
